\section{Discussion}

Our comprehensive evaluation framework reveals that the reflection mode delivers the optimal reliability-coverage balance among the three extraction modes. The reflection paradigm demonstrates superior performance across multiple complementary metrics: it achieves the highest CheckRules compliance (64.8\% for all four rules), generates the most triples per chunk (15.8), and wins a clear plurality in LLM-as-a-Judge evaluations across precision, comprehensiveness, and relevance dimensions.

The iterative critic-corrector loop systematically addresses schema violations while expanding triple coverage, yielding a denser yet cleaner knowledge graph. This improvement stems from the reflection mechanism's ability to identify and correct extraction errors while discovering previously missed but valid triples. The reflection mode's superior entity coverage ratio (ECR: 0.53 vs. 0.30-0.31) and relationship coverage ratio (RCR: 0.38 vs. 0.21-0.22) demonstrate its effectiveness in capturing diverse semantic content. The diversity analysis further illuminates the reflection mode's approach: while achieving the highest coverage, it exhibits lower entropy across all dimensions, indicating a deliberate reduction in uncertainty to yield a more compact, connected, and navigable graph.

However, these gains come with some trade-offs. The reflection agent requires additional inference rounds, potentially limiting its suitability for real-time applications requiring fast turnaround (e.g., intraday news feeds). In such scenarios, the single pass strategy may offer a viable alternative, recovering most normalization benefits with reduced computational overhead.

Our evaluation results reveal some limitations. First, cross document co-reference resolution is only partially addressed as the reflection loop operates on isolated filings. Second, our evaluation methodology depends on LLM-voting surrogate ground truth, risking propagation of biases inherent in the underlying judge models. 